\chapter{Análise Pós-Laboratorial}

A etapa de resultados consolida os dados numéricos extraídos das simulações transientes no LTSpice e os confronta com os valores nominais estipulados no projeto teórico (calculados via Maxima). O projeto do conversor \textit{Buck} estabeleceu como requisitos fundamentais uma tensão de saída contínua $V_0 = 2,5\text{ V}$ (metade da tensão de alimentação $V_{CC} = 5\text{ V}$) e uma ondulação máxima de tensão ($\Delta V$) estritamente limitada a $0,1\%$ de $V_{CC}$, o que corresponde a $5,0\text{ mV}$.

As Tabelas \ref{tab:resultados_v0} e \ref{tab:resultados_deltav} apresentam a comparação direta entre as grandezas teóricas e as medições experimentais obtidas com os cursores do simulador para as quatro combinações de carga e frequência.

\begin{table}[H]
	\centering
	\caption{Comparação dos valores de Tensão Média de Saída ($V_0$).}
	\label{tab:resultados_v0}
	\renewcommand{\arraystretch}{1.3}
	\begin{tabular}{cccccc}
		\hline
		\textbf{Cenário} & \textbf{Resistência ($R$)} & \textbf{Indutância ($L$)} & \textbf{$V_0$ Teórico} & \textbf{$V_0$ Simulado} & \textbf{Erro (\%)} \\ \hline
		1 & $5\,\Omega$ & $125\text{ mH}$ & $2,50\text{ V}$ & $2,9477\text{ V}$ & $17,91\%$ \\
		2 & $5\,\Omega$ & $25\text{ mH}$ & $2,50\text{ V}$ & $2,5016\text{ V}$ & $0,06\%$ \\
		3 & $1\,\Omega$ & $25\text{ mH}$ & $2,50\text{ V}$ & $2,4957\text{ V}$ & $0,17\%$ \\
		4 & $1\,\Omega$ & $5\text{ mH}$ & $2,50\text{ V}$ & $2,4976\text{ V}$ & $0,09\%$ \\ \hline
	\end{tabular}
\end{table}

\begin{table}[H]
	\centering
	\caption{Comparação dos valores de Ondulação de Tensão ($\Delta V$).}
	\label{tab:resultados_deltav}
	\renewcommand{\arraystretch}{1.3}
	\begin{tabular}{cccccc}
		\hline
		\textbf{Cenário} & \textbf{Resistência ($R$)} & \textbf{Indutância ($L$)} & \textbf{$\Delta V$ Teórico} & \textbf{$\Delta V$ Simulado} & \textbf{Erro (\%)} \\ \hline
		1 & $5\,\Omega$ & $125\text{ mH}$ & $5,00\text{ mV}$ & $4,9511\text{ mV}$ & $0,98\%$ \\
		2 & $5\,\Omega$ & $25\text{ mH}$ & $5,00\text{ mV}$ & $4,8635\text{ mV}$ & $2,73\%$ \\
		3 & $1\,\Omega$ & $25\text{ mH}$ & $5,00\text{ mV}$ & $4,9165\text{ mV}$ & $1,67\%$ \\
		4 & $1\,\Omega$ & $5\text{ mH}$ & $5,00\text{ mV}$ & $4,8600\text{ mV}$ & $2,80\%$ \\ \hline
	\end{tabular}
\end{table}
